\section{Limitations and Future Work}\label{sec:limitations}

This section acknowledges the open limitations of the Bailout Protocol as currently specified, describes the conditions under which each limitation is most acute, and identifies three productive directions for future research and engineering work that would strengthen the protocol's applicability in larger-scale and higher-stakes deployments.

\subsection{Limitations}\label{sec:limitations:subsection}

The following limitations are not implementation defects correctable by refinement within the current specification. They are structural properties of the protocol's design that impose real constraints on its applicability.

\subsubsection{Human Response Latency}\label{sec:human-response-latency}

The Bailout Protocol's liveness guarantee (Theorem~2) depends on the assumption that the Purpose Layer human anchor will engage with the bailout signal and issue a resolution within a time interval compatible with the operational requirements of the system. In slow-cycle domains---regulatory compliance review, strategic planning, multi-party contract negotiation---this assumption is well-grounded. In real-time cyber-physical domains, it is not.

Consider an autonomous agricultural system performing sub-second irrigation adjustments in response to microclimate sensor data. If the Bounds Engine fires a bailout because it encounters an unresolvable constraint state (for example, a water rights conflict between co-located parcels under different ownership), the protocol requires the human anchor to review the complete diagnostic context and issue a resolution before the Fact Layer will release the hard block on actuators. The human anchor may beunavailable, traveling, or cognitively engaged in an unrelated task. The $T_{\text{human}}$ timeout (2 hours, as specified in Section~\ref{sec:timeout-handling}) is architecturally conservative precisely because it does not impose a deadline on human judgment; but that conservativeness means the system may be stalled for a duration that is unacceptable in time-sensitive operational contexts.

The limitation is fundamental: the protocol trades real-time performance for accountability guarantee. For any deployment in which human review latency exceeds the operational time budget, the protocol either must be supplemented with a safe fallback state that preserves the hard block but maintains minimum viable operation, or the protocol's applicability must be scoped away from that operational tier. Neither workaround is built into the current specification.

\subsubsection{Adversarial Conditions}\label{sec:adversarial-conditions}

The protocol is designed to operate correctly under conditions of component degradation, environmental novelty, and operator error. It is not designed to operate correctly under conditions of active adversarial interference with the escalation path itself.

An adversarial agent or external actor with access to the system hierarchy could, in principle, interfere with bailout signal delivery by severing or corrupting the parent pointer chain, by denying network connectivity to the human anchor's notification channel, or by impersonating the human anchor to receive and resolve a bailout with fraudulent authority. The protocol's forensic ledger provides retrospective detection of such interference (the \texttt{HIERARCHY\_ERROR} and \texttt{BYPASS\_VIOLATION} markers are logged when propagation fails or override attempts are detected), but it does not prevent interference in real time.

The practical implication is that the protocol's security guarantees hold under the assumption that the system hierarchy and communication channels are trusted infrastructure. For deployment in adversarial network environments (for example, multi-agent systems operating across organizational trust boundaries without a shared identity and access management infrastructure), additional integrity mechanisms---authenticated parent chains, signed escalation messages, multi-factor human anchor verification---would be required before the protocol's guarantees hold. These mechanisms are not specified in the current design.

\subsubsection{Scalability}\label{sec:scalability}

The Bailout Protocol's correctness guarantees are proved for a system topology consisting of a single Bounds Engine node firing a single BailoutTrigger upward along a linear parent chain to a single human anchor. This model is adequate for small-scale deployments and for formal analysis. It is not adequate for multi-agent systems with dozens or hundreds of concurrent nodes, each capable of firing bailouts independently.

Under high node density, the following failure modes emerge:

\begin{itemize}\setlength\itemsep{0.25em}
\item \textbf{Anchor saturation.} A single human anchor assigned as \texttt{owner()} to $n$ concurrent nodes may receive $n$ simultaneous BailoutTriggers, each requiring genuine review. Human cognitive bandwidth is not a scalable resource. The protocol currently provides no mechanism for distributing anchor assignments across a pool of qualified human reviewers, nor any mechanism for priority-based ordering of concurrent bailout signals.
\item \textbf{Hierarchy breadth.} Deep linear hierarchies (the model assumed in correctness proofs) and wide flat hierarchies (more common in practice) exhibit different bailout propagation dynamics. In a wide hierarchy, a single bailout at a leaf node propagates to a designated human anchor whose authority over the broader system may be partial or ambiguous. The correctness proofs do not cover this case.
\item \textbf{Cascading bailout storms.} Under system-wide stress conditions (network partition, sensor degradation, or coordinated environmental shock), multiple nodes may fire bailouts simultaneously. The protocol does not specify any behavior for coordinating concurrent bailout signals at the same or different anchor nodes, nor any graceful degradation policy when bailout volume exceeds human review capacity.
\end{itemize}

These scalability failure modes are known and acknowledged. They do not invalidate the protocol's correctness guarantees under its stated assumptions; they define the boundaries of those assumptions for deployment planning.

\subsubsection{Exploitation Risk}\label{sec:exploitation-risk}

The Bailout Protocol creates a single, well-defined choke point through which all unresolvable exception states must pass: the Purpose Layer human anchor. A malicious actor who can trigger the protocol at will---by crafting inputs that satisfy trigger condition T1 (constraint uncovered) or T5 (scope exceeded) repeatedly---can subject the human anchor to a denial-of-service pattern in which their attention is consumed by a high volume of low-stakes bailout signals, rendering them unavailable for genuine high-stakes exceptions.

This exploitation risk is distinct from the adversarial conditions described above. In the adversarial case, the attacker targets the escalation infrastructure itself. In the exploitation case, the attacker targets the human anchor's attention, using the protocol's mandatory escalation requirement as the attack vector. The protocol's \texttt{SCOPE\_EXCEEDED} trigger condition (T5) is particularly exposed: an attacker with the ability to submit inputs outside the authorized scope of a node can fire a bailout at will for any node whose bounds are insufficiently narrow.

The current specification does not include any mechanism for rate-limiting bailout signals, for distinguishing genuine scope-exceeding events from adversarial input manipulation, or for automated pre-screening of bailout signals before they reach the human anchor. These are significant operational security gaps that must be addressed before deployment in adversarial environments.

\subsection{Future Work}\label{sec:future-work}

Three directions for future research and engineering are identified as the most productive near-term extensions of the current work.

\subsubsection{Formal Verification}\label{sec:formal-verification}

The correctness properties established in Section~\ref{sec:correctness-properties} are stated as proof sketches with explicit assumptions. A full formal verification of the Bailout Protocol state machine using a mechanized proof assistant (Coq, Isabelle, or TLA$+$) would transform these sketches into machine-checked theorems.

The verification effort would target three specific claims: (1) that the state machine transition function $\delta$ has no undefined transitions for any valid state-event pair in $\Sigma \times E$; (2) that the escalation algorithm terminates in at most $d$ steps for any node at depth $d$ in a finite rooted hierarchy; and (3) that the hard block $\Gamma_B$ is monotonic with respect to the state ordering, meaning the block is never released before the \texttt{RESOLVED} state is reached.

Mechanized verification is particularly important for the bypass mechanism (Section~\ref{sec:bypass-mechanism}), where the claim that no machine actor may resolve a trigger at any state other than \texttt{RESOLVED} is load-bearing for the entire accountability architecture. A mechanized proof that this property holds for all valid state transitions would close the principal remaining gap between the informal correctness argument and a formal guarantee.

Additionally, formal verification of the protocol's composition with the Recursive Spawning Protocol (RSP) would establish that the upstream accountability guarantee holds across nested spawn chains, not only for single-level nodes. The RSP introduces parent chain depth as a variable that affects the distance a BailoutTrigger must travel to reach a human anchor, and the interaction of this depth variable with the timeout constants $T_{\text{bound}}$ and $T_{\text{human}}$ requires careful analysis.

\subsubsection{Adaptive Timeout}\label{sec:adaptive-timeout}

The current specification uses fixed timeout constants $T_{\text{bound}}=300$\,s and $T_{\text{human}}=7200$\,s. These constants were chosen conservatively to accommodate thorough human review under normal operating conditions. They are not adapted to the operational context of the deployment.

An adaptive timeout extension would adjust timeout values based on three signals: (1) the operational criticality of the affected node (as defined by the Purpose Layer owner at provisioning); (2) the current backlog of pending bailout signals at the designated human anchor; and (3) historical data on human response times for comparable trigger condition classes. For a node operating in a non-critical decorative context, longer timeouts with lighter staleness warnings may be appropriate. For a node operating in a safety-critical context, shorter timeouts with escalation to a secondary anchor may be required.

The adaptive timeout mechanism must preserve the protocol's core invariant: no autonomous resolution is ever permitted regardless of timeout behavior. Adaptation should affect only the warning cadence and secondary anchor escalation timing, not the conditions under which the \texttt{RESOLVED} state may be entered. A productive formal framework for this extension is a timeout policy function $\theta : \mathbb{K} \times \mathbb{C} \times \mathbb{N} \rightarrow \mathbb{R}^+$ where $\mathbb{K}$ is the trigger condition class, $\mathbb{C}$ is the node criticality tier, and $\mathbb{N}$ is the current anchor backlog, producing a context-appropriate timeout value.

\subsubsection{Distributed Human Anchor Pools}\label{sec:anchor-pools}

The scalability limitations described in Section~\ref{sec:scalability} could be substantially mitigated by replacing the single-anchor model with a distributed human anchor pool: a defined set of qualified human reviewers, any one of whom may receive and resolve a bailout signal, with an allocation policy that distributes signals across the pool based on workload, expertise, and availability.

The distributed anchor pool extension requires three additional components. First, an \textbf{anchor discovery protocol} that resolves the question of which human anchors are available and qualified at the time a bailout fires, given that the Purpose Layer owner designated at provisioning may be unavailable. Second, an \textbf{allocation policy} that selects among available anchors using criteria that preserve the single-accountable-decision-maker property: exactly one anchor must receive each trigger, and that anchor must have genuine authority over the affected node's mandate. Third, a \textbf{distribution audit mechanism} that records in the Forensic Ledger which anchor received each bailout, enabling retrospective accountability reconstruction in cases where the allocation decision is questioned.

The extension is architecturally compatible with the current protocol: the escalation algorithm (Algorithm~\ref{alg:escalation}) resolves a single parent node per iteration, and the distributed pool extension replaces the static $\text{owner}(B)$ designation with a dynamic lookup against the pool registry. The state machine transitions and the hard block mechanism are unchanged. This makes the distributed anchor pool a natural layer-3 extension rather than a protocol redesign.
